\documentclass[sigconf,nonacm]{acmart}

\settopmatter{printacmref=false}
\renewcommand\footnotetextcopyrightpermission[1]{}
\pagestyle{plain}

\usepackage{booktabs}

\title{Open, Reproducible Consumption-Based Carbon Intensity Forecasting with Explicit Cross-Border Flow Modeling}

\author{Nicol\'as Higuera Wilches}
\affiliation{
  \institution{IE University, School of Science and Technology}
  \city{Madrid}
  \country{Spain}
}
\email{nicolaswilches@icloud.com}

\author{Bissan Ghaddar}
\affiliation{
  \institution{Advisor, IE University}
  \city{Madrid}
  \country{Spain}
}

\begin{document}

\begin{abstract}
We present an open, reproducible deep learning framework for short-horizon
consumption-based carbon intensity forecasting that explicitly models
cross-border electricity flows, evaluated across five structurally diverse grid
zones (US-MIDA-PJM, US-NY-NYIS, Finland, Belgium, and Singapore). The framework
reproduces the closed CarbonCast methodology as a production-based baseline,
extends it to a consumption-based target with interconnector flow and partner
carbon intensity inputs, and compares both against a single-tier model that
forecasts carbon intensity directly. We report a per-region accuracy atlas, a
structural ablation of the two-tier design, and a characterization of where
forecasts succeed and fail. % TODO: fold in final headline numbers once the
% Friday model comparison completes.
\end{abstract}

\maketitle

\section{Introduction and Problem}
% Reuse and tighten the proposal's framing.
Data center electricity demand is rising sharply, and the carbon caused by a
workload depends on when and where it runs. The operationally correct signal for
a carbon-aware scheduler is \emph{consumption-based} carbon intensity, which
accounts for the carbon embedded in imported electricity, rather than the
\emph{production-based} intensity of locally generated power that most academic
forecasters target. This report develops short-horizon (96 hour) forecasts of
consumption-based carbon intensity, opens a closed methodology, and characterizes
performance across heterogeneous grids.

\section{Related Work}
% TODO: CarbonCast (Maji et al., BuildSys '22), EnsembleCI, CarbonFlex, Phoenix.
% Position the gap: closed operational models; production-based academic targets;
% flows ignored. See docs/lit-notes.

\section{Data}
We use Electricity Maps (academic license) for hourly carbon intensity, power
breakdown, and electricity flows, and Open-Meteo (ERA5 reanalysis for history,
GFS for forecasts) for weather, over 2021 through May 2026 across the five zones.
Carbon intensity is computed with CarbonCast Table~1 lifecycle emission factors.
For each zone we also collect the carbon intensity history of its interconnector
partners (20 external zones), used as a consumption-based feature.

\paragraph{A regime shift in Finland.}
Annual mean consumption-based carbon intensity in Finland falls from about
161~gCO\textsubscript{2}eq/kWh in 2021 to about 80 in 2023, then stabilizes near
66 to 84 thereafter, while renewable share stays flat. The drop is consistent
with new nuclear capacity entering the mix. Belgium, by contrast, is stable over
the same window (roughly 150 to 195). This motivates a per-zone training window
(Section~\ref{sec:setup}): training Finland on the pre-shift years would teach a
regime that no longer holds.

\section{Methodology}
\subsection{Two-tier architecture}
Following CarbonCast, the two-tier model couples per-source Tier~1 feedforward
networks that forecast generation, with a Tier~2 convolutional-recurrent network
that consumes those forecasts together with historical carbon intensity, weather,
and calendar features to emit a 96 hour carbon intensity vector. We extend Tier~1
with per-interconnector flow networks that forecast signed net flow, and we add
partner carbon intensity as a Tier~2 channel. To avoid in-sample leakage, Tier~1
forecasts used for Tier~2 training are generated out of fold (leave one year out);
final Tier~1 models drive inference.

\subsection{Production and consumption targets}
The production-based variant (the faithful baseline) targets the domestic
generation mix intensity. The consumption-based variant targets the
Electricity Maps reported consumption intensity and adds the cross-border
channels.

\subsection{Single-tier baseline}
To test whether the two-tier structure earns its complexity, we compare against a
single-tier model: the same Tier~2 network run standalone, with no Tier~1 forecast
channels, predicting carbon intensity directly from its history plus weather and
calendar. Because the network is identical, the comparison isolates the
contribution of the two-tier design.

\section{Experimental Setup}
\label{sec:setup}
\paragraph{Splits.} Training 2021 through 2025, validation January through April
2026, Test~A May 2026, and a live Test~B (June 8 to 21 2026) head-to-head against
the Electricity Maps operational forecast. Finland uses a shortened training start
of 2023 (per-zone window), for the regime-shift reason above; the other zones use
the full window.

\paragraph{Training.} RMSE loss, Adam with gradient clipping. Each model is
trained on three random seeds and the model with the best validation score is
kept (Section~\ref{sec:findings}).

\section{Results}
% TODO: fill from outputs/friday_comparison.csv once runs complete.
Table~\ref{tab:atlas} reports Test~A accuracy per zone for the consumption-based
two-tier model. % Placeholder values from the full-window run; Finland to be
% updated with the shortened window.

\begin{table}[t]
\caption{Test A (May 2026) per-zone accuracy, consumption-based two-tier
(validation-selected). Provisional; see TODO.}
\label{tab:atlas}
\begin{tabular}{lrrr}
\toprule
Zone & MAPE (\%) & MAE & RMSE \\
\midrule
Singapore       & 0.90  & 4.17  & 6.54  \\
US-NY-NYIS      & 7.03  & 19.22 & 25.45 \\
US-MIDA-PJM     & 8.30  & 31.69 & 40.25 \\
Finland         & 24.04 & 11.43 & 14.41 \\
Belgium         & 64.42 & 74.17 & 90.23 \\
\bottomrule
\end{tabular}
\end{table}

% TODO: single-tier vs two-tier comparison table (the structural ablation).
% TODO: per-horizon degradation curve. TODO: forecast-vs-actual time series.

\section{Iterative Findings and Modeling Decisions}
\label{sec:findings}
Our methodology was shaped by a sequence of empirical findings, which we record
here because they justify the final design.

\paragraph{Initialization sensitivity.} Early five-zone runs showed large
run-to-run variation driven purely by the random initial weights: on the harder
zones, some seeds converged well while others diverged (to non-finite loss or to
a poor solution). On Belgium the seed-to-seed spread in percentage error exceeded
the size of any plausible modeling improvement.

\paragraph{Seeding and validation-based selection.} We therefore fixed the random
seed for reproducibility and adopted a selection rule: train several seeds and
keep the one with the best \emph{validation} score, never the test score. This
removes the seed lottery, because a divergent seed also scores poorly on
validation and is excluded automatically.

\paragraph{Gradient clipping.} The non-finite-loss failures were exploding
gradients; adding gradient clipping removed them. Stable seeds are unaffected.

\paragraph{Belgium as an outlier.} Belgium showed both the highest error and the
largest instability. Even after stabilization, its best-validation model
generalized poorly from the validation months to Test~A (validation near 33,
Test~A near 64), indicating a genuine distribution shift in May 2026 rather than a
training artifact. This is reported as a failure mode, not corrected away.

\paragraph{Finland regime shift and shortened training.} Inspecting carbon
intensity over time revealed the Finnish regime shift described in Section~3. We
shortened Finland's training window to begin in 2023 and re-ran the models.
% TODO: report the effect of the shortened window on Finland's metrics.

\section{Discussion and Limitations}
The Electricity Maps academic forecast horizon is 24 hours rather than 72, so the
head-to-head comparison (Test~B) is restricted to the 1 to 24 hour band.
% TODO: Belgium failure-mode analysis; single vs separate model structures.

\section{Conclusion}
% TODO: synthesize once the comparison is complete.

\end{document}
